
\begin{tabular}{|L{2.5cm}|C{1.8cm}|L{10.5cm}|}
	\hline
	\textbf{Period} & \textbf{Actor} & \textbf{Milestone / Contribution} \\
	\hline
	1960s & NASA (Apollo programme) & Introduced physical "twin" concept: an identical ground replica of the spacecraft used to simulate and troubleshoot the orbiting vehicle in real time. First practical predecessor of the DT. \\
	\hline
	1991 & David Gelernter & Published \textit{Mirror Worlds}, envisioned software models that mirror reality from information received from the physical world. Conceptual forerunner to DT. \\
	\hline
	2002 & Michael Grieves (Univ. of Michigan) & Formalised the Mirrored Spaces Model within PLM context: three components: Real Space, Virtual Space, and a bidirectional data/information link. First academic formulation of DT. \\
	\hline
	2006 & Grieves (revision) & Renamed model to Information Mirroring Model; emphasised bidirectionality and multiple virtual spaces for a single physical asset. \\
	\hline
	2010 & NASA (Glaessgen \& Stargel) & Coined the term "Digital Twin" in NASA's Technology Roadmap. Defined DT as "an integrated multi-physics, multi-scale, probabilistic simulation of a vehicle or system that uses the best available physical models, sensor updates, and fleet history to mirror the life of its flying twin". \\
	\hline
	2014 & Grieves (White Paper) & Published the three-component DT model: (a) physical products in Real Space, (b) virtual products in Virtual Space, (c) connections of data and information between real and virtual. \\
	\hline
	2016–2018 & Tao, Zhang et al. (Beihang Univ.) & Proposed the five-dimension DT model: Physical Entity (PE), Virtual Entity (VE), Services (Ss), Twin Data (DD), and Connections (CN). Now the dominant reference architecture. \\
	\hline
	2018 & Kritzinger et al. (TU Wien) & Published the most widely cited DT taxonomy (>736 citations): distinguished Digital Model, Digital Shadow, and Digital Twin based on data-flow automation and directionality. \\
	\hline
	2020–present & Global Industry \& Academia & DT adopted across aerospace, manufacturing, healthcare, smart cities, energy, and additive manufacturing. \\
	\hline
\end{tabular}